\bibitem{ROGStevenson}
G.~Stevenson, \emph{Subcategories of singularity categories via tensor actions},
\href{https://arxiv.org/abs/1105.4698v3}{arXiv:1105.4698v3},
Theorem~6.13 (PDF p.~26),
\href{https://doi.org/10.1112/S0010437X1300746X}{doi:10.1112/S0010437X1300746X}.
原論文の非次数付き分類を引用し、次数シフト軌道への移行は本文で証明する。
\bibitem{ROGDyckerhoff}
T.~Dyckerhoff, \emph{Compact generators in categories of matrix factorizations},
Duke Math. J. \textbf{159} (2011), 223--274;
\href{https://arxiv.org/abs/0904.4713v5}{arXiv:0904.4713v5},
\S2.4, \S5.6, Theorem~5.9. 本文の三次Euler計算は直接も証明する。
\bibitem{ROGStacks}
The Stacks Project Authors, \emph{The Stacks Project},
\href{https://stacks.math.columbia.edu/tag/09SM}{Tag 09SM},
\S13.37, Lemmas~13.37.3--4 and Proposition~13.37.6;
\href{https://stacks.math.columbia.edu/tag/00FV}{Tag 00FV}, Theorem~10.34.1;
\href{https://stacks.math.columbia.edu/tag/0BQQ}{Tag 0BQQ}, \S36.16 (projective subscheme generators).
\bibitem{ROGTV}
B.~To\"en and M.~Vaqui\'e, \emph{Moduli of objects in dg-categories},
\href{https://arxiv.org/abs/math/0503269v5}{arXiv:math/0503269v5},
Lemmas~2.8 and 2.10 (PDF pp.~21--24).
片側perfect性の比較と、perfect対象・射のfiltered-colimit降下に用いる。
